(1) 小球纯滚动约束为自转角速度 \(\omega=v/r\)。由高度 h 处静止到圆环上极角 θ（自底部向上测量）时球心升高 \(H=R(1-\cos\theta)\)，机械能守恒：
\[
mgh=\frac12mv^2+\frac12I(v/r)^2+mgR(1-\cos\theta)
=\frac12m(1+\lambda)v^2+mgR(1-\cos\theta).
\]
所以
\[
v^2(\theta)=\frac{2g[h-R(1-\cos\theta)]}{1+\lambda}.
\]

(2) 对球心沿半径指向轨道圆心方向取正，重力沿内向的分量为 \(-mg\cos\theta\)（底部指向圆心为向上），故牛顿向心方程
\[
N-mg\cos\theta=m v^2/R,\quad
N(\theta)=m[v^2(\theta)/R+g\cos\theta].
\]
在顶部 θ=π，支持力 \(N_{\rm top}=m[v_{\rm top}^2/R-g]\)，因此不能只要求顶部速度非零。

(3) 完整接触需要所有位置 N≥0，最低约束来自顶部；它要求 \(v_{\rm top}^2\ge gR\)。顶部高度为 2R，能量给
\[
v_{\rm top}^2=\frac{2g(h-2R)}{1+\lambda}.
\]
联立得到
\[
h\ge h_{\min}=2R+\frac{1+\lambda}{2}R=\frac{5+\lambda}{2}R.
\]
为证明顶部条件确为全局最严格，代入速度公式：
\(N/(mg)=2(h/R-1+\cos\theta)/(1+\lambda)+\cos\theta\)，这是 cosθ 的严格递增线性函数，故在顶部 cosθ=−1 取最小值。

(4) 当 λ=0，\(h_{\min}=5R/2\)，恢复无自转点质点在竖直圆环的临界释放高度。均匀实心球 λ=2/5 时
\[
h_{\min}=\frac{27}{10}R.
\]
增加的 \(R/5\) 来自滚动旋转动能，因此小球实际滚动比无转动点质点需要更高释放高度。真实轨道还要求连接段几何上可实现给定无滑动初态。
